On éclaire une fente verticale de largeur $a=\qty{5e-5}{\m}$, à l'aide d'un
Laser qui donne une lumière monochromatique de longueur d'onde
$\lambda=\qty{633}{\nm}$. On observe sur un écran placé à la distance
$D$ de la fente, des taches lumineuses. La mesure de la largeur
de la tache centrale à donné $L=\qty{3.8}{\cm}$.

La célérité de la lumière dans le vide (ou l'air) est $c=\qty{3e8}{\v}$.

\begin{minipage}{0.59\linewidth}
    \flushleft
    \begin{tikzpicture}[x={(190:7mm)},y={(-10:1cm)},z={(90:1cm)}]
        \def\L{6}
        \def\H{3}
        \def\l{1}
        \def\h{1.5}
        \def\a{0.15}
        \def\A{2}
        \def\B{0.8}
        \def\C{0.6}
        \def\D{5.5}
        \def\d{8}
        % Écran 1
        \begin{scope}[canvas is yz plane at x=0]
            \path[draw,fill=black!5] (-\L/2,-\H/2) rectangle (\L/2,\H/2);
            \node[above left,rotate=-10] at (\L/2,\H/2) {Écran};
            \draw[->] (-0.55*\L,0) node[above] {$x^{\prime}$} --
                (0.55*\L,0) node[above] {$x$};
            \draw[->] (0,-0.6*\H) node[right] {$y^{\prime}$} --
                (0,0.6*\H) node[right] {$y$};
        \end{scope}
        % Écran 2
        \begin{scope}[canvas is yz plane at x=\D]
            \path[draw,fill=black!20] (-\l/2,-\h/2) rectangle (\l/2,\h/2);
            \path[draw,fill=white] (-\a/2,-0.4*\h) rectangle (\a/2,0.4*\h);
            \draw[<-] (-\a/2,0.3*\h) -- ++(160:2)
                node[above] {Fente de largeur $a$};
        \end{scope}
        % Rayon
        \draw[black!50,thick] (\d,0,0) -- ++(-\d+\D,0,0);
        % Laser
        \begin{scope}[canvas is xy plane at z=\C/2]
            \path[draw,fill=black!20] (\d,-\B/2) rectangle (\d+\A,\B/2);
        \end{scope}
        \begin{scope}[canvas is zx plane at y=\B/2]
            \path[draw,fill=black!20] (-\C/2,\d) rectangle (\C/2,\d+\A);
            \node[below=1mm] at (-\C/2,\d+\A/2) {Laser};
        \end{scope}
        \begin{scope}[canvas is yz plane at x=\d+\A]
            \path[draw,fill=black!20] (-\B/2,-\C/2) rectangle (\B/2,\C/2);
        \end{scope}
    \end{tikzpicture}
\end{minipage}
\hfill
\begin{minipage}{0.40\linewidth}
    \flushright
    \begin{tikzpicture}
        \newdimen\a \a=3mm
        \newdimen\D \D=4.5cm
        \newdimen\L \L=1.7cm
        \pgfmathparse{atan(\L/(2\D))}\let\th=\pgfmathresult
        \node[fill,inner sep=0pt,minimum width=1mm,minimum height=4cm] (D) {};
        \node[fill=white,inner sep=0pt,minimum width=3mm,
              minimum height=\a] (F) at (D.center) {};
        \draw[<->,>={Straight Barb[round,length=3pt,width=3pt]}]
              (F.south west) -- (F.north west) node[pos=0.5,left] {$a$};
        \node[ellipse,inner sep=0pt,fill=black!30,minimum height=\L,
              minimum width=4mm] (T0) at ([xshift=\D]F.east) {};
        \foreach \y [count=\i] in {1.5cm,-1.5cm}
        \node[ellipse,inner sep=0pt,fill=black!30,minimum height=0.6\L,
              minimum width=3mm] (T\i) at ([xshift=\D,yshift=\y]F.east) {};
        \draw[densely dashed] (F.east) -- ($(T0.north)!0.5!(T1.south)$);
        \draw[densely dashed] (F.east) -- (T0.center);
        \draw[->] ([xshift=3cm]F.east) arc (0:\th:3) node[pos=0.5,right] {$\theta$};
        \draw[{Bar[] . Straight Barb[]}-{Straight Barb[] . Bar[]}]
            ([xshift=5mm]T0.north) -- ++(0,-\L) node[right,pos=0.5] {$L$};
        \draw[{Bar[] . Straight Barb[]}-{Straight Barb[] . Bar[]}]
            ([yshift=2mm]D.north east) -- ++(\D,0) node[above,pos=0.5] {$D$};
    \end{tikzpicture}
\end{minipage}
\begin{enumerate}
    \item Nommer le phénomène observé sur l'écran.~\textbf{(0,5pt)}
    \item Sur quelle direction ($xx^{\prime}$ ou $yy^{\prime}$) s’étalent les
          taches obtenues.~\textbf{(0,5pt)}
    \item Rappeler la relation qui lie  $\theta$, $\lambda$ et $a$.~\textbf{(0,5pt)}
    \item Montrer que l'écart angulaire $\theta$ s'exprime par~:
          $\theta=\frac{L}{2 \cdot D}$
          (on prend~: $\tan\theta \simeq \theta$ (\unit{\rad})).~\textbf{(0,75pt)}
    \item Établir l'expression de la distance $D$ en fonction de $L$, $a$ 
          et $\lambda$. Calculer $D$.~\textbf{(0,75pt)}
    \item Déterminer la valeur de la fréquence $f$ de la lumière utilisée.~\textbf{(1,0pt)}
    \item On refait l'expérience en utilisant un fil très fin vertical de
          diamètre $a_0$, on obtient une tache centrale de largeur
          $L_0=\qty{1.5}{\cm}$. Déterminer la valeur de $a_0$.~\textbf{(1,0pt)}
    \item On éclaire avec cette source Laser un prisme de verre flint d'indice
          de réfraction $n=1,64$. 
          \begin{enumerate}
              \item Déterminer la vitesse $v$ de propagation de la lumière dans
                    le prisme.~\textbf{(1,0pt)}
              \item Trouver la longueur d’onde $\lambda^{\prime}$ au cours de la
                    propagation dans le prisme.~\textbf{(1,0pt)}
              \item Qu'observe-t-on si on remplace la lumière monochromatique
                    par la lumière blanche~? Quel est le nom de ce phénomène~?~\textbf{(1,0pt)}
          \end{enumerate}
\end{enumerate}

